\NativeTitle{9}{第二单元素养评价（A）}{（相交与平行）}{40}{5}{一 & 二 & 三 & 四 & 五}{ &  &  &  & }
\NativeSlot{283}{76}{%
\ExamGroup{一、填空。{\normalfont\kaishu\fontsize{10.8bp}{14bp}\selectfont （每空1分，共19分）}}

% p009_s01

\ExamQuestion{1.}{5}%
下面的各组线中，\mbox{（\hspace{1.8em}）}互相平行；\mbox{（\hspace{1.8em}）}相交，\mbox{（\hspace{1.8em}）}互相垂直。（填序号）\par
\BookFigure{figures/p009_s01-01.tex}
}
\NativeSlot{418}{73}{%
% p009_s02

\ExamQuestion{2.}{5}%
下面哪个图形中，既有互相平行的线段，又有互相垂直的线段？在它下面画“\(\checkmark\)”。\par
\BookFigure{figures/p009_s02-01.tex}
}
\NativeSlot{545}{44}{%
% p009_s03

\ExamQuestion{3.}{5}%
同一平面内，两条直线\(m\)和\(n\)的位置关系只有\mbox{（\hspace{1.8em}）}和\mbox{（\hspace{1.8em}）}，如果直线\(m\)和\(n\)相交成直角，那么直线\(m\)和\(n\)互相\mbox{（\hspace{1.8em}）}，记作\mbox{（\hspace{1.8em}）}，读作\mbox{（\hspace{7em}）}。\par
}
\NativeSlot{621}{36}{%
% p009_s04

\ExamQuestion{4.}{5}%
从客厅的吊灯点到地面画三条线段，长度分别是3米、4米、5米，其中一条是垂线段，那么吊灯点到地面的距离是\mbox{（\hspace{1.8em}）}米。\par
}
\NativeSlot{672}{36}{%
% p009_s05

\ExamQuestion{5.}{5}%
长方形操场的长是100米，宽60米，操场的对边互相\mbox{（\hspace{1.8em}）}，邻边互相\mbox{（\hspace{1.8em}）}。沿着跑道跑一圈，跑过的路程是\mbox{（\hspace{1.8em}）}米。\par
}
\NativeSlot{722}{36}{%
% p009_s06

\ExamQuestion{6.}{5}%
把一张长方形纸沿同一方向对折两次，展开后的折痕互相\mbox{（\hspace{1.8em}）}，如果沿横竖不同方向对折两次，折痕互相\mbox{（\hspace{1.8em}）}。\par
}
\NativeSlot{774}{36}{%
% p009_s07

\ExamQuestion{7.}{5}%
平行的健身步道两边之间，能画\mbox{（\hspace{1.8em}）}条垂线段，这些线段的长度都\par
}
